\chapter{L'immunité et la réponse immunitaire}

\noindent\textit{Notre organisme est en permanence exposé à des microbes potentiellement
dangereux. Il dispose pourtant de moyens de défense efficaces, qui s'organisent en
plusieurs lignes successives. Ce chapitre étudie ces défenses, la réaction inflammatoire,
et le cas particulier des greffes.}
\vspace{8pt}

\connecteBanner

\section{L'immunité, une défense de l'organisme}

\begin{definition}[Immunité]
L'\textbf{immunité} est l'ensemble des mécanismes par lesquels l'organisme se
\textbf{défend} contre les agents pathogènes (bactéries, virus, champignons\dots) et les
corps étrangers.
\end{definition}

\begin{propriete}[Première ligne de défense : les barrières naturelles]
La \textbf{peau} et les \textbf{muqueuses} (yeux, bouche, voies respiratoires) forment une
barrière physique et chimique qui empêche la plupart des microbes de pénétrer dans
l'organisme. Lorsque cette barrière est franchie (plaie, piqûre), d'autres défenses entrent
en jeu.
\end{propriete}

\section{La réaction inflammatoire}

\begin{definition}[Réaction inflammatoire]
La \textbf{réaction inflammatoire} est une réponse locale et rapide de l'organisme face à
une infection ou une blessure. Elle se manifeste par quatre signes caractéristiques :
\textbf{rougeur}, \textbf{chaleur}, \textbf{douleur} et \textbf{gonflement}.
\end{definition}

\begin{propriete}[Mécanisme]
Au site de l'infection, les vaisseaux sanguins se dilatent (d'où la rougeur et la
chaleur) et deviennent plus perméables, ce qui permet à de nombreux \textbf{phagocytes}
(globules blancs) de sortir du sang et de converger vers la zone infectée.
\end{propriete}

\begin{definition}[Phagocytose]
La \textbf{phagocytose} est le mécanisme par lequel un phagocyte \textbf{englobe} puis
\textbf{détruit} un microbe.
\end{definition}

\begin{illus}
\begin{tikzpicture}[scale=0.85]
\draw[onbo,line width=1.3pt,fill=onbo!10] (0,0) circle(1.1);
\node[font=\small] at (0,0) {phagocyte};
\filldraw[onbo2,line width=0pt] (2.6,0) circle(0.2);
\node[font=\scriptsize] at (2.6,-0.5) {microbe};
\draw[->,onbmuted,line width=1pt,dashed] (1.9,0) -- (1.25,0);
\node[above,font=\scriptsize,onbmuted] at (1.6,0.15) {englobement};
\end{tikzpicture}
\fig{Figure — La phagocytose : le phagocyte englobe puis détruit le microbe.}
\end{illus}

\begin{remarque}
La fièvre, souvent associée à une infection, aide également l'organisme à se défendre :
une température plus élevée ralentit la multiplication de nombreux microbes.
\end{remarque}

\section{Deux formes d'immunité}

\begin{propriete}[Immunité naturelle et immunité acquise]
L'\textbf{immunité naturelle} (ou innée) est \textbf{présente dès la naissance},
\textbf{rapide} mais \textbf{non spécifique} : elle réagit de la même façon face à
n'importe quel agent pathogène (c'est le cas de la phagocytose). L'\textbf{immunité
acquise} (ou adaptative) se \textbf{met en place progressivement}, est \textbf{spécifique}
d'un agent pathogène précis, et laisse une \textbf{mémoire immunitaire} (étudiée au
chapitre suivant).
\end{propriete}

\section{Les greffes et le rejet}

\begin{definition}[Greffe]
Une \textbf{greffe} est le transfert d'un organe ou d'un tissu (le \textbf{greffon})
d'un \textbf{donneur} vers un \textbf{receveur}.
\end{definition}

\begin{propriete}[Trois types de greffes]
\begin{center}
\small
\begin{tabular}{lll}
\toprule
\textbf{Type} & \textbf{Donneur / receveur} & \textbf{Risque de rejet} \\
\midrule
Autogreffe & même individu & aucun \\
Allogreffe & même espèce, individus différents & possible \\
Xénogreffe & espèces différentes & très élevé \\
\bottomrule
\end{tabular}
\end{center}
\end{propriete}

\begin{propriete}[Le rejet de greffe]
Le système immunitaire du receveur peut reconnaître le greffon comme un corps
\textbf{étranger} et chercher à le \textbf{détruire} : c'est le \textbf{rejet}. Pour
limiter ce risque, on recherche une bonne \textbf{compatibilité} entre donneur et
receveur, et on prescrit des médicaments \textbf{immunosuppresseurs} qui atténuent la
réponse immunitaire du receveur.
\end{propriete}

\begin{exemple}
Une greffe de peau prélevée sur soi-même (autogreffe, par exemple après une brûlure) ne
provoque pas de rejet, contrairement à une greffe de rein provenant d'un autre individu
(allogreffe), qui nécessite un traitement immunosuppresseur à vie.
\end{exemple}

% ═══════════════════════════════════════════════════════════════
\section{L'essentiel à retenir}

\begin{synthese}
\textbf{Immunité.} Ensemble des défenses de l'organisme contre les agents pathogènes.\\[4pt]
\textbf{Réaction inflammatoire.} Rougeur, chaleur, douleur, gonflement ; afflux de
phagocytes.\\[4pt]
\textbf{Phagocytose.} Un phagocyte englobe puis détruit un microbe.\\[4pt]
\textbf{Immunité naturelle} (rapide, non spécifique) \textbf{vs immunité acquise}
(spécifique, avec mémoire).\\[4pt]
\textbf{Greffes.} Autogreffe (aucun rejet), allogreffe (rejet possible), xénogreffe
(rejet très élevé).\\[4pt]
\textbf{Piège fréquent.} Le rejet de greffe n'est pas un dysfonctionnement : c'est le
système immunitaire qui fonctionne normalement en reconnaissant un corps étranger.
\end{synthese}

% ═══════════════════════════════════════════════════════════════
\section{Méthodes \& savoir-faire}

\methode{Reconnaître les signes d'une réaction inflammatoire}
Rechercher les quatre signes caractéristiques : rougeur, chaleur, douleur, gonflement.
\begin{exemple}
Une zone rouge, chaude et gonflée autour d'une écharde correspond à une réaction
inflammatoire.
\end{exemple}

\methode{Distinguer immunité naturelle et immunité acquise}
Se demander si la réponse est immédiate et identique pour tout agent pathogène (naturelle)
ou spécifique et progressive (acquise).
\begin{exemple}
La phagocytose, qui agit contre tout microbe, relève de l'immunité naturelle.
\end{exemple}

\methode{Identifier un type de greffe et son risque de rejet}
Comparer le donneur et le receveur : même individu (autogreffe), même espèce mais
individus différents (allogreffe), espèces différentes (xénogreffe).
\begin{exemple}
Une greffe de cornée provenant d'un autre être humain est une allogreffe.
\end{exemple}

% ═══════════════════════════════════════════════════════════════
\section{Exercices d'application}
\begin{multicols}{2}

\rubrique{Immunité et barrières}
\exo{}\diff{1} Qu'est-ce que l'immunité ?

\exo{}\diff{2} Citer deux barrières naturelles qui protègent l'organisme des microbes.

\rubrique{Réaction inflammatoire}
\exo{}\diff{1} Citer les quatre signes caractéristiques d'une réaction inflammatoire.

\exo{}\diff{2} Pourquoi la zone infectée devient-elle rouge et chaude ?

\exo{}\diff{2} Qu'est-ce que la phagocytose ?

\rubrique{Deux formes d'immunité}
\exo{}\diff{2} Quelle est la différence entre immunité naturelle et immunité acquise ?

\exo{}\diff{2} Laquelle des deux formes d'immunité est spécifique d'un agent pathogène
précis ?

\rubrique{Greffes}
\exo{}\diff{1} Qu'est-ce qu'une greffe ?

\exo{}\diff{2} Quel type de greffe ne présente aucun risque de rejet ? Pourquoi ?

\exo{}\diff{2} Qu'est-ce qu'un médicament immunosuppresseur ?

% ═══════════════════════════════════════════════════════════════
\end{multicols}

\section{Exercices d'entraînement}
\begin{multicols}{2}

\rubrique{Immunité et barrières}
\exo{}\diff{3} Expliquer pourquoi une plaie ouverte augmente le risque d'infection.

\rubrique{Réaction inflammatoire}
\exo{}\diff{3} Expliquer pourquoi une réaction inflammatoire, bien que désagréable, est
utile à l'organisme.

\exo{}\diff{2} Pourquoi la fièvre accompagne-t-elle souvent une infection ?

\exo{}\diff{3} Un patient présente une zone rouge, chaude, douloureuse et gonflée autour
d'une plaie. Que se passe-t-il probablement à cet endroit ?

\rubrique{Deux formes d'immunité}
\exo{}\diff{3} Pourquoi dit-on que l'immunité naturelle est « non spécifique » ?

\exo{}\diff{3} Pourquoi l'immunité acquise est-elle plus lente à se mettre en place que
l'immunité naturelle ?

\rubrique{Greffes}
\exo{}\diff{2} Classer, du risque de rejet le plus faible au plus élevé : allogreffe,
autogreffe, xénogreffe.

\exo{}\diff{3} Pourquoi recherche-t-on une bonne compatibilité entre donneur et receveur
avant une allogreffe ?

\exo{}\diff{3} Un patient greffé doit prendre des immunosuppresseurs à vie. Quel est
l'inconvénient de ce traitement pour sa capacité à se défendre contre les infections ?

\exo{}\diff{2} Une greffe de peau prélevée sur la cuisse d'un patient et posée sur son
bras est-elle une autogreffe ou une allogreffe ? Justifier.

% ═══════════════════════════════════════════════════════════════
\end{multicols}

\section{Sujets type BEPC}

\sujetbac[La réaction inflammatoire]
\begin{enumerate}[label=\arabic*)]
\item Cite les quatre signes caractéristiques d'une réaction inflammatoire.
\item Qu'est-ce que la phagocytose ?
\item Explique pourquoi les vaisseaux sanguins se dilatent au site d'une infection.
\item En quoi cette réaction est-elle utile à l'organisme ?
\end{enumerate}

\sujetbac[Immunité naturelle et immunité acquise]
\begin{enumerate}[label=\arabic*)]
\item Définis l'immunité naturelle.
\item Définis l'immunité acquise.
\item Cite une caractéristique qui différencie ces deux formes d'immunité.
\item Donne un exemple de mécanisme relevant de l'immunité naturelle.
\end{enumerate}

\sujetbac[Greffes et rejet]
\begin{enumerate}[label=\arabic*)]
\item Définis les termes « greffe », « donneur » et « receveur ».
\item Cite les trois types de greffes et classe-les selon leur risque de rejet.
\item Qu'est-ce que le rejet de greffe ?
\item À quoi servent les médicaments immunosuppresseurs prescrits après une greffe ?
\end{enumerate}

% ═══════════════════════════════════════════════════════════════
\section{Corrigés}
\begin{multicols}{2}

\corrige{de l'exercice 1}
L'ensemble des mécanismes par lesquels l'organisme se défend contre les agents
pathogènes.

\corrige{de l'exercice 2}
La peau et les muqueuses.

\corrige{de l'exercice 3}
Rougeur, chaleur, douleur, gonflement.

\corrige{de l'exercice 4}
Car les vaisseaux sanguins s'y dilatent, augmentant l'afflux de sang.

\corrige{de l'exercice 5}
Le mécanisme par lequel un phagocyte englobe puis détruit un microbe.

\corrige{de l'exercice 6}
L'immunité naturelle est rapide et non spécifique ; l'immunité acquise est spécifique
d'un agent pathogène et laisse une mémoire immunitaire.

\corrige{de l'exercice 7}
L'immunité acquise.

\corrige{de l'exercice 8}
Le transfert d'un organe ou d'un tissu d'un donneur vers un receveur.

\corrige{de l'exercice 9}
L'autogreffe, car le greffon provient du receveur lui-même : le système immunitaire ne le
reconnaît pas comme étranger.

\corrige{de l'exercice 10}
Un médicament qui atténue la réponse immunitaire du receveur, pour limiter le risque de
rejet.

\corrige{du sujet 1}
1) Rougeur, chaleur, douleur, gonflement.\\
2) Le mécanisme par lequel un phagocyte englobe puis détruit un microbe.\\
3) Pour permettre un afflux plus important de sang et de phagocytes vers la zone
infectée.\\
4) Elle permet de mobiliser rapidement les défenses de l'organisme contre l'infection.

\corrige{du sujet 2}
1) Une immunité présente dès la naissance, rapide et non spécifique.\\
2) Une immunité qui se met en place progressivement, spécifique d'un agent pathogène, et
qui laisse une mémoire immunitaire.\\
3) Par exemple : la rapidité (immunité naturelle) contre la spécificité (immunité
acquise).\\
4) La phagocytose.

\corrige{du sujet 3}
1) Greffe : transfert d'un organe/tissu (greffon) d'un donneur vers un receveur.\\
2) Autogreffe (aucun risque), allogreffe (risque possible), xénogreffe (risque très
élevé).\\
3) La reconnaissance du greffon comme corps étranger par le système immunitaire du
receveur, qui cherche alors à le détruire.\\
4) Ils atténuent la réponse immunitaire du receveur pour limiter le risque de rejet du
greffon.

\end{multicols}
\exercicesBanner
